Augmenting GenUI with Design Patterns

% Promise: What is a technology that promises us something good?
Language models (LMs) can now generative high-fidelity UI screens (GenUI) based on a user's natural language description.

% Problem: What is the problem that prevents this technology from realizing its promise?
However, the GenUI process remains a ``black box'' that lacks transparency and control: it is often unclear how LMs come up with a particular UI design or how a designer can steer LMs to follow a specific direction beyond prompt engineering.

% Prior work: What is the most representative prior work that has attempted to solve this problem and why did it fail?
To address this limitation, prior work is mainly divided into two approaches:
\one Before generation---improving how users convey their design intent beyond text prompt [\xx] and;
\two After generation---enabling targeted iteration, such as decomposing a generated UI into slots wherein designers can have fine control over specific levels of elements [\xx].
None of these approaches directly taps into the actual generation process.

% Proposed solution: What is your proposed solution given where prior work fell short?
We propose a complementary approach---using interaction design patterns (hereafter referred to as ``patterns'' for brevity) as an intermediate layer to enhance the transparency and control of GenUI.
Specifically, a pattern is an abstract specification of a prototypical type of UI screen common amongst multiple applications. For example, for many e-commerce mobile apps (\eg shopping or food delivery), the pattern is a list of categories where each row horizontally presents items of that category.
Our approach first maps a prompt to a pattern, which we then use to condition LMs' generation of UI code, \eg using specific elements from a design system to instantiate a pattern.
We implemented this pattern-based approach as a model context provider (MCP), called PatternMCP, and demonstrate its usefulness on PatternGenUI--a front-end where a designer enters a prompt to generate a UI screen, selects and/or edit relevant patterns retrieved from a curated library (currently consisting of [N] patterns), and obtains a pattern-guided generated UI from an LM.

% Proof of contribution: What is your proof to validate that your proposed solution contributes to solving that problem?
To validate our approach, we conducted a two-fold evaluation:
First, an ablative evaluation shows that users often preferred UI screens generated with PatternMCP than the ones without it even though the underlying LM was the same.
Further, a designer evaluation identifies both strengths and trade-offs of PatternGenUI compared to a state-of-the-art professional tool. [more detailed findings on transparency and control]
